
Only 1 in 1,000 Python files in a curated code corpus contains an independently-executable
doctest --- a $31\times$ gap below the 3.1\% rate at which \texttt{>>>} doctest patterns appear. This finding
emerges from a systematic three-phase feasibility scan of 10,000 randomly sampled Python files:
Phase A (pattern detection: 3.1\%), Phase B (AST extraction: 2.0\%), and Phase C (subprocess
execution: 0.1\%). The collapse at Phase C is driven by import isolation: Python library code
depends on third-party packages unavailable in clean subprocess environments. The resulting
executable-doctest token pool (0.004M tokens estimated) is approximately 125,$000\times$ below a
500M-token supervised fine-tuning (SFT) budget target, making doctest-passing filtering
structurally infeasible for raw Python corpus SFT construction at this scale. We validate
a three-phase filtering pipeline (28/28 unit tests passing, 129.8 seconds for 10,000 files)
and identify compile()-based filtering as the practical execution-quality alternative ---
scalable to 12.96M+ files without dependency management. The equal-token-budget SFT
comparison experiment (compile-only vs. unfiltered, Qwen2.5-Coder-1.5B, HumanEval/MBPP
pass@1) is designed and designated as future work.

